% !PRINTED 508
\weberpage{508}
Namely, let $\mu, \nu, \ldots$ be the degrees of the equations with integral coefficients which (by 2.) the functionals $\alpha, \beta, \ldots$ satisfy, and $m=\mu\nu\ldots$ By $\omega_1, \omega_2, \ldots, \omega_m$ we understand the $m$ quantities
\[
\alpha^r\beta^s\ldots \qquad r=0,1,\ldots,\mu-1, \quad s=0,1,\ldots,\nu-1, \ldots
\]
Then, by means of the equations of degrees $\mu, \nu, \ldots$ which the numbers $\alpha, \beta, \ldots$ satisfy, all products $\alpha^r\beta^s\ldots$ in which one of the exponents $r, s, \ldots$ is greater than $\mu-1, \nu-1, \ldots$ can be expressed linearly in the form (5) in terms of $\omega_1, \omega_2, \ldots, \omega_m$, and the same therefore also holds of every functional $\omega$, and consequently also of the products $\omega\omega_1, \omega\omega_2, \ldots, \omega\omega_m$. From this it follows by 6. that $\omega$ is integral, as we wished to prove.

As a special case of the theorem 7., from which, incidentally, the general theorem can easily be inferred again, we emphasise:

8. \gesperrt{By addition, subtraction and multiplication of integral functionals there always arise integral functionals again.}

We have already remarked earlier (§.\,134) that one can always determine an algebraic field which contains arbitrarily given algebraic numbers. From this it follows that in the theorem 7. the $\alpha, \beta, \ldots$ may be arbitrary integral algebraic numbers or functionals, and the corresponding statement holds of the theorem 8.

From the definition 2. it follows further, by the theorem §.\,136, 8.:

9. \gesperrt{If $\omega$ is an integral functional of the field $\Omega$, then all functionals conjugate to $\omega$ are also integral.}

As a particular consequence of this theorem, the following may further be mentioned:

10. \gesperrt{The absolute norm of an integral functional is a natural integer.}

Finally we prove also the following theorem:

11. \gesperrt{If a functional $\omega$ satisfies an equation of the form}
\weq{10}{\omega^m + \alpha_1\omega^{m-1} + \cdots + \alpha_m = 0}
\gesperrt{in which $\alpha_1, \alpha_2, \ldots, \alpha_m$ are integral functionals in $\Omega$, then $\omega$ is also an integral functional.}
